% !TeX root = ../main.tex

\chapter{Conclusion and Outlook}\label{ch:conclusion}

This chapter summarises the answers to the four research questions posed in
\autoref{sec:research-questions}, restates the scope those answers carry, and states the
work this thesis leaves open.

\section{Concluding Remarks}\label{sec:concluding-remarks}

RQ1 asked whether longitudinal modelling adds information beyond static \ac{FC}
representations. Under the evaluated protocol, neither age and sex alone nor a single
visit's connectivity alone separates converters from stable \ac{MCI} above chance, while a
linear model reading only the change in connectivity between two visits already reaches
within 0.04 \ac{AUC} of the fully trained temporal-first configuration selected in
\autoref{sec:rq3}. The change between visits, rather than any single visit's state,
therefore carries essentially all of the classification signal available under this
protocol. A subject-count-matched sensitivity check found that restricting a trained
model's own predictions to subjects with three or more visits raises its score, while
training a separate model on only that restricted subset lowers it; the loss from a
smaller training pool outweighs the gain from longer visit histories at this cohort size,
which establishes only that the two effects could not be separated here, not that longer
visit histories carry no additional signal.

RQ2 asked whether spatial-first graph representation learning adds value. On the
\ac{DELCODE}-only protocol, a configuration with no graph encoder at all scores within
fold-to-fold spread of a frozen, reconstruction-pretrained encoder while using
approximately seventeen times fewer trainable parameters, and the pretrained encoder's own
reconstruction fidelity falls in the band this thesis labels poor. Reconstruction
pretraining does outperform training the same encoder from a random initialisation end to
end, which is consistent with pretraining stabilising a difficult optimisation problem
rather than supplying a representation that improves on the absence of any learned
encoder. Fine-tuning the pretrained encoder jointly with the classification head introduces
a further, bit-reproducible instability across seeds that three candidate mechanisms could
explain, none of which is isolated by an experiment run for this thesis. Under the
evaluated protocol, spatial-first graph representation learning does not add value over
mean-pooled connectivity features at this cohort size.

RQ3 asked whether temporal-first modelling improves over spatial-first modelling. On the
pooled protocol, the selected temporal-first configuration (S1) reaches a pooled out-of-fold
\ac{AUC} of $0.7488 \pm 0.0033$, exceeding the randomly initialised spatial-first floor by
0.186 \ac{AUC} at fourteen times fewer trainable parameters and the frozen pretrained
spatial-first floor by 0.030 \ac{AUC} at 7.6 times fewer trainable parameters. Importantly,
S1 contains neither a graph autoencoder nor graph message-passing layers; its empirical
advantage stems entirely from preserving region-resolved temporal trajectory information before
spatial pooling, rather than from graph-based spatial aggregation. Two configurations that
add parameters on top of the selected configuration, a temporal saliency gate and a
dual-score head, both score lower than the configuration they extend, so this pattern is not
explained by the selected configuration simply carrying more capacity than its comparators.
Under the evaluated \ac{ADNI}-and-\ac{DELCODE} protocol, temporal-first modelling performs well
relative to every spatial-first comparator evaluated, including the published Brain-TokenGT
competitor, which under the evaluated configuration showed substantially greater run-to-run
variability and scored below every temporal-first configuration that clears the pre-specified floor
gates of \autoref{sec:protocol}.

RQ4 asked whether these findings transfer to an independent cohort. The selected
configuration and its two pre-specified secondaries all reach in-domain test \ac{AUC} in
the high 0.7s, but cross-cohort transfer to \ac{OASIS} was at chance under the evaluated
protocol, with every configuration's four-seed mean falling within one standard error of an
\ac{AUC} of 0.5. A linear probe recovers cohort of origin from the pooled representation at
approximately 0.86 \ac{AUC}, well above the pre-specified escalation threshold of 0.75,
despite no cohort label ever entering the training objective; this strong decodability is
consistent with substantial cohort dependence in the learned representation and may
contribute to the absence of transfer, though no experiment in this thesis isolates that
causal path. A pre-specified gradient-reversal mitigation, applied once this threshold was
found exceeded, moved both the cohort-decodability score and pooled out-of-fold \ac{AUC} in
the direction opposite to its design intent, so this particular mitigation, as configured,
did not resolve the dependence it targeted.

These four answers are scoped to the protocol under which they were obtained: an in-domain
comparison of longitudinal architectures on a pooled \ac{ADNI}-and-\ac{DELCODE} cohort of
248 subjects, evaluated out-of-fold across four seeds, together with a single one-shot read
of a held-out in-domain test set and \ac{OASIS} for the selected configuration and its
pre-specified secondaries. They do not license a claim of external validation, a
time-to-event prediction, or a healthy-reference or individualized-deviation framework,
none of which this protocol was designed to test.

\section{Future Work}\label{sec:future-work}

Developing cohort-invariant models based on harmonised clinical target definitions and
individualized deviation from an appropriate reference population remains a topic for
future work.

Isolating the contribution of the inter-visit interval requires a cohort whose visit
spacing is irregular enough to exercise it: over ninety percent of \ac{DELCODE}'s recorded
intervals fall at twelve months (\autoref{sec:comparability}), so this contribution remains
untested by the present evaluation, and a cohort with more irregular spacing, such as
\ac{ADNI} or \ac{OASIS}, would be needed to test it directly. Whether the graph
construction stage itself, rather than only the encoder that consumes it, contributes
information beyond mean-pooled connectivity is a related open question the encoder
ablation of \autoref{sec:rq2} does not answer, since it varies encoder initialisation
without varying how the graph is built in the first place. Replacing the reconstruction
pretext objective with one less prone to spending representational capacity on
subject-specific idiosyncrasy, as the account in \autoref{sec:ssl-contribution} suggests,
is a direct response to that account rather than a demonstration that any specific
alternative objective would perform better.

Recovering full statistical power for the interpretability analyses of
\autoref{sec:interpretability-results} at a larger cohort size, and running the
parcellation ablation left untested at the fixed 200-region resolution used throughout
\autoref{sec:problem-formulation}, are each direct extensions of a specific gap this thesis
identifies but does not close. Adding a time-to-event head that predicts the interval to
conversion rather than only its binary outcome, an extension the protocol of this thesis
explicitly does not license (\autoref{sec:scope-of-claims}), would require a survival model
fit to the same cohorts under the same leakage controls used throughout
\autoref{ch:methodology}.

Fusing the functional-connectivity representation evaluated here with structural or
fluid-biomarker measures was outside the scope of the present single-modality evaluation
and remains open. Testing a domain-adaptation mitigation against the cross-cohort transfer
question of \autoref{sec:cohort-dependence} directly, rather than only diagnosing the
representation's cohort dependence as this thesis does, is a further extension. Finally,
adopting same-seed replication as a reporting protocol, rather than as an audit applied
after the fact to one published competitor, would let a future evaluation at this cohort
size distinguish an architecture's own instability from ordinary run-to-run noise before a
comparison is drawn, a distinction the reproducibility audit of \autoref{sec:limitations}
shows cannot otherwise be assumed.
